% !TEX TS-program = pdflatexmk
\RequirePackage{iftex}
\ifPDFTeX
  \pdfoutput=1
  \pdfminorversion=7
  \input glyphtounicode
  \pdfgentounicode=1
\fi
\ifdefined\pdfinterwordspaceon\pdfinterwordspaceon\fi
\documentclass[justified,notoc,twoside,nobib,nofonts]{tufte-handout}
\usepackage[T1]{fontenc}
\usepackage[utf8]{inputenc}
\usepackage[english]{babel}
\ifPDFTeX
  \usepackage[tracking]{microtype}
\else
  \usepackage{microtype}
\fi
\usepackage{amsthm,thmtools,thm-restate}
\usepackage{amsmath,amssymb,mathtools}
% Text and mathematics.
\usepackage[nomath,oldstylenums]{kpfonts}
\edef\KPbodyfamily{\rmdefault}
\renewcommand{\rmdefault}{zpltlf}
\usepackage{newpxmath}
\let\rmdefault\KPbodyfamily
\renewcommand{\familydefault}{\rmdefault}

% BaskervilleF headings use native Type 1 fonts and work with pdfLaTeX.
\newcommand{\BaskervilleHeading}{%
  \fontencoding{T1}\fontfamily{BaskervilleF-LF}%
  \fontseries{m}\fontshape{n}\selectfont}
\titleformat{\section}[hang]
  {\BaskervilleHeading\fontsize{15}{18}\selectfont}
  {\thesection}{1em}{}
\titleformat{\subsection}[hang]
  {\BaskervilleHeading\fontsize{13}{16}\selectfont}
  {\thesubsection}{1em}{}
\titleformat{\paragraph}[runin]
  {\BaskervilleHeading}{\theparagraph}{1em}{}
\usepackage{etoolbox}
\makeatletter
\patchcmd{\maketitle}{\LARGE\itshape\@title}
  {\BaskervilleHeading\fontsize{22}{26}\selectfont\raggedright\@title}
  {}{\PackageError{two-circuit-four-matchings}{Title font patch failed}{}}
% The author and email are printed immediately below the title block.
\patchcmd{\maketitle}{\@@author\@@date}{\@@date}
  {}{\PackageError{two-circuit-four-matchings}{Author block patch failed}{}}
\makeatother
\usepackage{booktabs,array,enumitem,needspace}
\usepackage[numbers,square]{natbib}
\usepackage{xcolor}
\usepackage{xurl,hyperref,cleveref}
\urlstyle{same}
\geometry{a4paper,top=23mm,bottom=46mm,left=28mm,right=56mm,marginparwidth=3cm}
\setlength{\emergencystretch}{2em}
\clubpenalty=10000
\widowpenalty=10000
\displaywidowpenalty=10000
\setcounter{secnumdepth}{2}
\AtBeginDocument{\fontsize{12}{15}\selectfont}
\fancyhead[LE]{\thepage\quad\textsmallcaps{nikolay ulyanov}}
\fancyhead[RO]{\textsmallcaps{four perfect matchings}\quad\thepage}
\declaretheorem[name=Theorem,numberwithin=section]{theorem}
\declaretheorem[name=Proposition,sibling=theorem]{proposition}
\declaretheorem[name=Lemma,sibling=theorem]{lemma}
\declaretheorem[name=Corollary,sibling=theorem]{corollary}
\declaretheorem[name=Definition,sibling=theorem,style=definition]{definition}
\setlist[enumerate]{label=\textup{(\roman*)},leftmargin=*,itemsep=2pt,topsep=4pt}
\crefname{theorem}{theorem}{theorems}
\crefname{lemma}{lemma}{lemmas}
\crefname{definition}{definition}{definitions}
\setcitestyle{numbers,square}

% Full references appear in the margin at their first citation.
% Keep the canonical margin-selection block below when cloning this paper.
\makeatletter
\newdimen\@gp@fnmarkwidth
\newcommand{\bothcite}[2][0cm]{%
  \citep{#2}\sidenote[][#1]{\raggedright\csname paper@#2\endcsname}%
  \kern\@gp@fnmarkwidth\@gp@cite@space}
\def\@gp@cite@space{\@ifnextchar.{\@gp@cite@halfspace}{\@gp@cite@comma}}
\def\@gp@cite@comma{\@ifnextchar,{\@gp@cite@halfspace}{\@gp@cite@semicolon}}
\def\@gp@cite@semicolon{\@ifnextchar;{\@gp@cite@halfspace}{\@gp@cite@colon}}
\def\@gp@cite@colon{\@ifnextchar:{\@gp@cite@halfspace}{\@gp@cite@exclamation}}
\def\@gp@cite@exclamation{\@ifnextchar!{\@gp@cite@halfspace}{\@gp@cite@question}}
\def\@gp@cite@question{\@ifnextchar?{\@gp@cite@halfspace}{\@gp@cite@openparen}}
\def\@gp@cite@openparen{\@ifnextchar({\@gp@cite@halfspace}{\@gp@cite@closeparen}}
\def\@gp@cite@closeparen{\@ifnextchar){\@gp@cite@halfspace}{\@gp@cite@bracket}}
\def\@gp@cite@bracket{\@ifnextchar]{\@gp@cite@halfspace}{\@gp@cite@normalspace}}
\def\@gp@cite@halfspace{\nobreak\hskip.5\fontdimen2\font\relax}
\def\@gp@cite@normalspace{\nobreak\space}
% GP margin selection v1 (pdfLaTeX). Keep this block in the standalone source.
\makeatletter
\@ifundefined{@gp@fnmarkwidth}{\newdimen\@gp@fnmarkwidth}{}
\ifPDFTeX
\def\@makefnmark{%
  \hbox{%
    \setbox\z@=\hbox{\@textsuperscript{\normalfont\footnotesize\@thefnmark}}%
    \global\@gp@fnmarkwidth=\wd\z@%
    \setbox\tw@=\copy\z@%
    \immediate\pdfxform\tw@%
    \pdfannot width\wd\z@ height\ht\z@ depth\dp\z@{%
      /Subtype /Watermark /F 708 /Contents ()%
      /AP << /N \the\pdflastxform\space 0 R >>%
    }%
  }%
}
\fi
% Selection-safe margin text: annotation appearances, with separate URI links.
\ifPDFTeX
\newbox\@gp@marginbox
\newdimen\@gp@marginwidth
\newdimen\@gp@marginheight
\newdimen\@gp@marginfirstheight
\newdimen\@gp@marginpadding
\newdimen\@gp@marginskip
\newdimen\@gp@margindepth
\newcount\@gp@marginserial
\long\def\@gp@marginappearance#1{%
  \begingroup
  \global\let\@gp@marginuri\@empty
  \setbox\@gp@marginbox=\vtop{%
    \hsize=\marginparwidth
    \def\href##1##2{%
      \ifx\@gp@marginuri\@empty
        \xdef\@gp@marginuri{\detokenize{##1}}%
      \else
        \PackageError{gp-margin-selection}{More than one URI in a margin note}%
          {Use a separate note for each reference.}%
      \fi
      \textcolor{\@urlcolor}{##2}%
    }%
    #1\par
  }%
  \@gp@marginfirstheight=\ht\@gp@marginbox
  \setbox\@gp@marginbox=\vbox{\unvbox\@gp@marginbox}%
  \@gp@marginpadding=\baselineskip
  \advance\@gp@marginpadding-\prevdepth
  \advance\@gp@marginpadding-\@gp@marginfirstheight
  \ifdim\@gp@marginpadding<\lineskiplimit
    \@gp@marginpadding=\lineskip
  \fi
  \@gp@marginskip=\baselineskip
  \advance\@gp@marginskip-\prevdepth
  \advance\@gp@marginskip-\ht\@gp@marginbox
  \ifdim\@gp@marginskip<\lineskiplimit
    \@gp@marginskip=\lineskip
  \fi
  \advance\@gp@marginpadding-\@gp@marginskip
  \vskip\@gp@marginpadding
  \@gp@marginwidth=\wd\@gp@marginbox
  \@gp@marginheight=\ht\@gp@marginbox
  \@gp@margindepth=\dp\@gp@marginbox
  \immediate\pdfxform\@gp@marginbox
  \edef\@gp@marginform{\the\pdflastxform}%
  \global\advance\@gp@marginserial\@ne
  \hbox{%
    \pdfannot width\@gp@marginwidth height\@gp@marginheight depth\@gp@margindepth{%
      /Subtype /Watermark /F 708 /Contents ()%
      /NM (GPmargin-\the\@gp@marginserial)%
      /AP << /N \@gp@marginform\space 0 R >>%
    }%
    \ifx\@gp@marginuri\@empty\else
      \pdfannot width\@gp@marginwidth height\@gp@marginheight depth\@gp@margindepth{%
        /Subtype /Link /F 4 /Border [0 0 0]%
        /A << /S /URI /URI (\pdfescapestring{\@gp@marginuri}) >>%
      }%
    \fi
    \vrule width0pt height\@gp@marginheight depth\@gp@margindepth
    \kern\@gp@marginwidth
  }%
  \endgroup
}
\else
\long\def\@gp@marginappearance#1{#1}
\fi
\renewcommand\@footnotetext[2][0pt]{%
  \marginpar{%
    \hbox{}\vspace*{#1}%
    \def\baselinestretch{\setspace@singlespace}%
    \reset@font\footnotesize%
    \@tufte@margin@par%
    \vspace*{-1\baselineskip}%
    \@gp@marginappearance{%
      \noindent
      \protected@edef\@currentlabel{%
        \csname p@footnote\endcsname\@thefnmark%
      }%
      \color@begingroup
        \@makefntext{\ignorespaces#2}%
      \color@endgroup
    }%
  }%
}
\renewcommand\marginnote[2][0pt]{%
  \ifthenelse{\boolean{@tufte@loadnatbib}}{%
    \let\cite\@tufte@infootnote@cite%
  }{}%
  \gdef\@tufte@citations{}%
  \marginpar{%
    \hbox{}\vspace*{#1}%
    \@tufte@marginnote@font\@tufte@marginnote@justification%
    \@tufte@margin@par\vspace*{-1\baselineskip}%
    \@gp@marginappearance{\noindent #2}%
  }%
  \@tufte@print@citations%
  \ifthenelse{\boolean{@tufte@loadnatbib}}{%
    \let\cite\@tufte@normal@cite%
  }{}%
}
\makeatother

\expandafter\def\csname paper@KM26\endcsname{%
  J\'an Karab\'a\v{s} and Edita M\'a\v{c}ajov\'a.
  \emph{On $4$-perfect-matching covers of cubic graphs with two adjacent
  odd circuits in a $2$-factor}.
  \href{https://arxiv.org/abs/2605.09475}{\nolinkurl{arXiv:2605.09475}}
  (2026).}
\expandafter\def\csname paper@LCKM18\endcsname{%
  C.~L. Lucchesi, M.~H. de Carvalho, N.~Kothari and U.~S.~R. Murty.
  \emph{On two unsolved problems concerning matching covered graphs}.
  SIAM Journal on Discrete Mathematics 32(2) (2018), 1478--1504.
  \href{https://doi.org/10.1137/17M1138704}{\nolinkurl{doi:10.1137/17M1138704}}.}
\expandafter\def\csname paper@Tutte47\endcsname{%
  W.~T. Tutte.
  \emph{The factorization of linear graphs}.
  Journal of the London Mathematical Society 22(2) (1947), 107--111.
  \href{https://doi.org/10.1112/jlms/s1-22.2.107}{\nolinkurl{doi:10.1112/jlms/s1-22.2.107}}.}
\expandafter\def\csname paper@KCLL20\endcsname{%
  N.~Kothari, M.~H. de Carvalho, C.~L. Lucchesi and C.~H.~C. Little.
  \emph{On essentially $4$-edge-connected cubic bricks}.
  Electronic Journal of Combinatorics 27(1) (2020), P1.22.
  \href{https://doi.org/10.37236/8594}{\nolinkurl{doi:10.37236/8594}}.}
\expandafter\def\csname paper@Edmonds65\endcsname{%
  Jack Edmonds.
  \emph{Maximum matching and a polyhedron with $0,1$-vertices}.
  Journal of Research of the National Bureau of Standards,
  Section B 69B (1965), 125--130.
  \href{https://doi.org/10.6028/jres.069B.013}{\nolinkurl{doi:10.6028/jres.069B.013}}.}
\expandafter\def\csname paper@CL00\endcsname{%
  C.~N. Campos and C.~L. Lucchesi.
  \emph{On the relation between the Petersen graph and the characteristic
  of separating cuts in matching covered graphs}.
  Technical Report IC-00-22, Institute of Computing,
  University of Campinas (2000).
  \href{https://www.ic.unicamp.br/~reltech/2000/00-22.pdf}{report}.}
\newcommand{\paperreference}[1]{\csname paper@#1\endcsname}
\makeatother

\newcommand{\Petersen}{P_{10}}
\newcommand{\three}{\{1,2,3\}}
\DeclareMathOperator{\odd}{odd}

\title[Graph Puzzles IV.2-preview: Four perfect matchings]{Graph Puzzles IV.2-preview: Four perfect matchings for snarks with a two-circuit $2$-factor}
\author{Nikolay Ulyanov + AI-Slop}
\date{6 October 2026}
\hypersetup{
  pdftitle={Graph Puzzles IV.2-preview: Four perfect matchings for snarks with a two-circuit 2-factor},
  pdfauthor={Nikolay Ulyanov + AI-Slop},
  colorlinks=true,
  linkcolor=blue!50!black,
  citecolor=blue!50!black,
  urlcolor=blue!50!black
}

\begin{document}
\selectlanguage{english}
\maketitle
Nikolay Ulyanov\textsuperscript{*} + AI-Slop%
\marginnote{\textsuperscript{*}{\scriptsize\href{mailto:ulyanick@gmail.com}{ulyanick@gmail.com}}}\par

\begin{abstract}
\noindent\newthought{Abstract: }
Let $G$ be a snark distinct from the Petersen graph, and let $M$ be a
perfect matching such that $G-M$ consists of two circuits. We prove that
$M$ belongs to a cover of $E(G)$ by four perfect matchings. Consequently,
$G$ has perfect matching index four. The proof combines the
Campos--Lucchesi theorem on separating cuts with the
Karab\'a\v{s}--M\'a\v{c}ajov\'a theorem on Hamiltonian cubic $3$-poles.
We show that the cut between the two factor circuits is separating,
use a perfect matching meeting this cut in three edges to reduce the
graph, and extend a cover of the reduced graph along odd paths.
The result applies, in particular, to every permutation snark distinct
from the Petersen graph.
\end{abstract}

\section{Introduction}

The \emph{perfect matching index} $\tau(G)$ of a bridgeless cubic graph
$G$ is the least number of perfect matchings whose union is $E(G)$.
Three perfect matchings cover $G$ precisely when they form the colour
classes of a $3$-edge-colouring. Covers by four perfect matchings are
therefore the next case for snarks.

Graphs are finite and loopless. Throughout this paper, a \emph{snark}
is a connected, bridgeless, simple
cubic graph of girth at least five and cyclic edge-connectivity at least
four that has chromatic index four. A graph is cyclically
$4$-edge-connected if each edge-cut whose two shores contain circuits
has size at least four. We write $\Petersen$ for the Petersen graph.
A \emph{circuit} is a connected $2$-regular subgraph, and a \emph{$2$-factor}
is a spanning $2$-regular subgraph.

Karab\'a\v{s} and M\'a\v{c}ajov\'a \bothcite{KM26} proved that a cubic
graph has a cover by four perfect matchings when it has a $2$-factor
consisting of two odd circuits joined by exactly three edges of the
complementary perfect matching. Their Conjecture~3.9 asks whether every
cubic graph with a two-circuit $2$-factor has perfect matching index at
most four, with the Petersen graph as the exception. We prove this
assertion for snarks, with an additional conclusion concerning the
complementary perfect matching.

\Needspace{9\baselineskip}
\begin{restatable}{theorem}{fourmatchings}\label{thm:main}
Let $G\not\cong\Petersen$ be a snark, and let $M$ be a perfect matching
such that $G-M=A\cup B$ has exactly two circuit components.
There are perfect matchings $M_1,M_2,M_3$ of $G$ such that
\[
  E(G)=M\cup M_1\cup M_2\cup M_3.
\]
In particular, $\tau(G)=4$.
\end{restatable}

An edge of $M$ joining $A$ and $B$ is a \emph{spoke}; one with both
ends on one circuit is a \emph{chord}. The theorem allows any number of
chords. A snark with a $2$-factor consisting of two induced circuits is
called a \emph{permutation snark}. We obtain the following consequence.

\begin{corollary}\label{cor:permutation}
Every permutation snark distinct from $\Petersen$ has perfect matching
index four. For each $2$-factor consisting of two induced circuits, its
complementary perfect matching belongs to a cover by four perfect
matchings.
\end{corollary}
\begin{proof}
Apply \cref{thm:main} to the complementary perfect matching of the
given $2$-factor.
\end{proof}

The cut between the two factor circuits is separating. A theorem of
Campos and Lucchesi gives a perfect matching meeting it in three edges.
We suppress odd paths to obtain a graph covered by the three-spoke
theorem, then extend its cover to $G$ while retaining $M$.

\section{Matching-covered contractions}

For $X\subseteq V(G)$, the cut $\delta_G(X)$ consists of the edges with
exactly one end in $X$; the sets $X$ and $V(G)\setminus X$ are its
\emph{shores}. A cut is \emph{trivial} if one shore consists of a single
vertex. We write $G/X$ for the graph obtained by contracting $X$ to
one vertex and deleting the resulting loops. Contraction and suppression
may produce parallel edges, which we retain as distinct edges. For an
edge set $S$, let $V(S)$ be the set of its end vertices. Path length
always counts edges.

We use the terminology of matching-covered graphs as in Lucchesi,
de Carvalho, Kothari and Murty \bothcite{LCKM18}.
A graph is \emph{matching covered} if it is connected and each edge
belongs to a perfect matching. A cut is \emph{separating} if both
cut-contractions are matching covered. It is \emph{tight} if every
perfect matching meets it in exactly one edge. A \emph{brick} is a
nonbipartite matching-covered graph whose tight cuts are all trivial.
A graph $J$ is \emph{factor-critical} if $J-v$ has a perfect matching
for every $v\in V(J)$.

\Needspace{9\baselineskip}
The following lemma will establish that the cut between $A$ and $B$
is separating.

\begin{lemma}\label{lem:contraction}
Let $H$ be a connected bridgeless loopless multigraph, and let
$b\in V(H)$. Suppose that $b$ has odd degree, every vertex of $H-b$
has degree three in $H$, and $H-b$ is factor-critical.
Then $H$ is matching covered.
\end{lemma}
\begin{proof}
Each edge $bu$ belongs to a perfect matching: take a perfect matching
of $H-b-u$ and add $bu$. Consider an edge $uv$ with $u,v\ne b$.
Suppose that $H-u-v$ has no perfect matching. By Tutte's theorem
\bothcite{Tutte47}, there is a set
$X\subseteq V(H)\setminus\{u,v\}$ such that $H-X-u-v$ has more than
$|X|$ odd components. Since $H-u-v$ has even order, the number $q$
of these components satisfies
\begin{equation}\label{eq:tutte}
  q\ge |X|+2.
\end{equation}

If $b\in X$, choose a perfect matching of $H-b-u$. Delete the $|X|$
vertices in $(X\setminus\{b\})\cup\{v\}$ from this graph.
Each odd component of the remaining graph contains a vertex whose
partner in the chosen matching was deleted. Distinct odd components
require distinct deleted partners. Hence $q\le |X|$, contradicting
\eqref{eq:tutte}.

If $b\notin X$, put $R=X\cup\{u,v\}$. All vertices of $R$ have degree
three, and $uv\in E(H[R])$, so
\begin{equation}\label{eq:upper}
  |\delta_H(R)|=3|R|-2|E(H[R])|\le 3|X|+4.
\end{equation}
All vertices of $H$ have odd degree. Each odd component $K$ of $H-R$
therefore has an odd edge boundary. Connectedness and bridgelessness
give $|\delta_H(V(K))|\ge 3$. These boundaries are disjoint subsets of
$\delta_H(R)$, and \eqref{eq:tutte} yields
\[
  |\delta_H(R)|\ge 3q\ge 3|X|+6,
\]
contrary to \eqref{eq:upper}. Thus $H-u-v$ has a perfect matching.
Adding $uv$ proves that this edge also belongs to a perfect matching.
\end{proof}

We also recall the standard connection between snarks and bricks.
For cubic matching-covered graphs, every tight cut has size three;
see Kothari, de Carvalho, Lucchesi and Little \bothcite[-2.5cm]{KCLL20},
Theorem~8 and Corollary~9. The argument below gives the needed consequence.

\begin{lemma}\label{lem:brick}
Every snark is a simple brick.
\end{lemma}
\begin{proof}
Let $G$ be a snark, and assign $x_e=1/3$ to every $e\in E(G)$.
This vector satisfies the degree equations
$x(\delta_G(v))=1$ for all $v\in V(G)$. For every odd vertex set $X$,
the cut $\delta_G(X)$ has odd size. Bridgelessness gives
$|\delta_G(X)|\ge 3$, and hence $x(\delta_G(X))\ge 1$.
Edmonds' perfect matching polytope theorem \bothcite[-3cm]{Edmonds65}
therefore expresses $x$ as a convex combination of incidence vectors
of perfect matchings. Since $x_e>0$ for every edge, $G$ is matching covered.

If $D$ is a tight cut, every perfect matching has exactly one edge
in $D$. The same convex combination gives
\[
  1=x(D)=\frac{|D|}{3},
\]
so $|D|=3$. Cyclic $4$-edge-connectivity implies that one shore of $D$
induces a forest. If that forest has $h$ vertices and $c$ components,
then its boundary has size
\[
  3h-2(h-c)=h+2c.
\]
Since $h,c\ge 1$ and $h+2c=3$, the shore consists of a single vertex.
Thus every tight cut is trivial. Finally, a cubic bipartite graph is
$3$-edge-colourable, so $G$ is nonbipartite and hence a brick.
\end{proof}

We use the following consequence of Campos and Lucchesi
\bothcite{CL00}, Theorem~1.1. Their result is stated for near-bricks;
the case of simple bricks is also recorded in
\citep{LCKM18}, footnote~2 on p.~1496.

\begin{theorem}[Campos--Lucchesi]\label{thm:CL}
Let $G\not\cong\Petersen$ be a simple brick, and let $C$ be a nontrivial
separating cut of $G$. There is a perfect matching $N$ of $G$ such that
\[
  |N\cap C|=3.
\]
\end{theorem}

We now apply \cref{lem:contraction} to the given $2$-factor.

\begin{lemma}\label{lem:separating}
Let $G$ be a snark, and let $M$ be a perfect matching for which
$G-M=A\cup B$ has two circuit components.
Then $A$ and $B$ are odd circuits, and
$C=\delta_G(V(A))$ is a nontrivial separating cut.
\end{lemma}
\begin{proof}
The circuits $A$ and $B$ have the same parity, because $G$ has even
order. If both were even, their edges could be coloured alternately
with two colours and the edges of $M$ with a third. Hence both are odd.
All edges of $C$ belong to $M$, so
\[
  |C|\equiv |V(A)|\equiv 1\pmod 2.
\]
The cut separates two circuits. Cyclic $4$-edge-connectivity therefore
gives $|C|\ge 5$.

Consider $H=G/V(B)$, with contraction vertex $b$.
Every vertex of $H-b$ has degree three in $H$, while
$d_H(b)=|C|$ is odd. The graph $H-b$ contains the spanning odd circuit
$A$, so deleting any vertex leaves a spanning path of even order and
therefore a perfect matching. Thus $H-b$ is factor-critical.
Each edge of $A$ lies on that circuit, each chord lies on a circuit
formed with a path of $A$, and each edge incident with $b$ lies on a
circuit formed with another such edge and a path of $A$. Consequently,
$H$ is connected and bridgeless. By \cref{lem:contraction}, it is
matching covered.

Interchanging $A$ and $B$ proves that $G/V(A)$ is matching covered as
well. Thus $C$ is separating. Both shores contain circuits of length
at least five, so $C$ is nontrivial.
\end{proof}

\section{The case of three spokes}

A cubic \emph{$3$-pole} has three dangling edges, each with a single
end vertex; every vertex is incident with three edges. A perfect
matching of a $3$-pole is a set of ordinary and dangling edges meeting
each vertex exactly once. If a $3$-pole $H$ has a spanning circuit $R$,
then $Q=E(H)\setminus E(R)$ is a perfect matching and contains all
three dangling edges.

Theorem~3.1 of Karab\'a\v{s} and M\'a\v{c}ajov\'a \citep{KM26} has
the following form.

\begin{theorem}[Karab\'a\v{s}--M\'a\v{c}ajov\'a]\label{thm:KM}
Let $H$ be a Hamiltonian cubic $3$-pole with a distinguished spanning
circuit $R$. There is a cover of $H$ by four perfect matchings, one
of which is $Q=E(H)\setminus E(R)$, such that every dangling edge
belongs to exactly two of the four matchings.
\end{theorem}

For the reduction in the next section, we need the conclusion with
the complementary matching specified. It follows by gluing the two
$3$-poles, as in \citep{KM26}, Corollary~3.7.

\begin{lemma}\label{lem:three-spokes}
Let $H$ be a loopless cubic multigraph, and let $S$ be a perfect matching
such that $H-S$ consists of two odd circuits. If exactly three edges
of $S$ join the two circuits, there are perfect matchings $L_1,L_2,L_3$
of $H$ such that
\[
  E(H)=S\cup L_1\cup L_2\cup L_3.
\]
\end{lemma}
\begin{proof}
Cut the three edges joining the circuits. Each circuit, together with
its chords and three dangling edges, forms a Hamiltonian cubic
$3$-pole. Apply \cref{thm:KM} to each pole, using label $4$ for the
matching complementary to its spanning circuit and labels $1,2,3$
for the other three matchings.

Each pole has odd order. A perfect matching of such a pole contains
an odd number of dangling edges. In the cover supplied by
\cref{thm:KM}, the three dangling edges occur once each among the
matchings labelled $1,2,3$. Each of these three matchings therefore
contains exactly one dangling edge. Permute the labels $1,2,3$ on
one pole so that the labels agree on the pairs of dangling edges
that are to be joined. Gluing gives four perfect matchings of $H$.
The matching labelled $4$ is precisely $S$, and every edge is covered.
\end{proof}

\Needspace{11\baselineskip}
\section{Proof of the main theorem}

\Needspace{8\baselineskip}
\fourmatchings*
\begin{proof}
By \cref{lem:brick}, $G$ is a simple brick. By
\cref{lem:separating}, the cut $C=\delta_G(V(A))$ is nontrivial and
separating. Theorem~\ref{thm:CL} supplies a perfect matching $N$ with
$|N\cap C|=3$. Put
\[
  S=M\cap N.
\]
Since $C\subseteq M$, exactly three edges of $S$ join $A$ and $B$.
The edges of $N\setminus M$ form a perfect matching of
\[
  (A\cup B)-V(S).
\]
Between consecutive vertices of $V(S)$ on either factor circuit,
the internal vertices therefore occur in even number. Every path
of the circuit between consecutive retained vertices has odd length.

Delete the edges of $M\setminus S$, and suppress the resulting vertices
of degree two along $A$ and $B$. Denote the resulting cubic multigraph
by $G'$. Its retained matching is $S$, and $G'-S=A'\cup B'$ consists
of two circuits. If $c_A$ and $c_B$ are the numbers of chords of $S$
on $A$ and $B$, respectively, then
\[
  |V(A')|=3+2c_A,\qquad |V(B')|=3+2c_B.
\]
Thus both circuits are odd and have at least three vertices. The graph
$G'$ is loopless, and exactly three edges of $S$ join $A'$ and $B'$.
By \cref{lem:three-spokes}, it has a cover
\[
  E(G')=S\cup L_1\cup L_2\cup L_3
\]
by four perfect matchings.

We extend this cover to $G$ using the alternating assignment along
odd paths from \citep{KM26}, Lemma~3.2. Give an edge $e$ of $G'$ the
set of labels
\[
  \lambda(e)=\{i\in\three:e\in L_i\}
       \cup\begin{cases}\{4\},&e\in S,\\ \varnothing,&e\notin S.\end{cases}
\]
\Needspace{6\baselineskip}
At each vertex, the label sets on the three incident edges partition
$\{1,2,3,4\}$. Each set is nonempty. If $e\in E(A'\cup B')$, then
$T=\lambda(e)$ is a subset of $\three$ of size one or two, because
the edge of $S$ at either end carries label $4$ and the other circuit
edge has a nonempty label set.

Replace $e$ by its original odd path in $A\cup B$. Give the successive
edges of that path the label sets
\begin{equation}\label{eq:extension}
  T,\quad \three\setminus T,\quad T,\quad
  \three\setminus T,\quad\ldots,\quad T.
\end{equation}
Keep the labels on every edge of $S$, and give each restored edge of
$M\setminus S$ the label set $\{4\}$.
At a retained vertex, the incident label sets agree with those in
$G'$. At a restored vertex, the two circuit edges carry complementary
subsets of $\three$, and the matching edge carries $\{4\}$.
Hence each of the four labels occurs on exactly one edge incident
with every vertex of $G$, and every edge carries at least one label.

For $i\in\three$, let $M_i$ be the set of edges carrying label $i$.
These are perfect matchings of $G$, and the edges carrying label $4$
are exactly the edges of $M$. Thus
\[
  E(G)=M\cup M_1\cup M_2\cup M_3.
\]
A cover of a cubic graph by three perfect matchings assigns each of
its three incident edges to a different matching at every vertex,
and therefore gives a $3$-edge-colouring. Since $G$ is a snark,
$\tau(G)=4$.
\end{proof}

\clearpage
\begin{thebibliography}{Edmonds65}
\bibitem[KM26]{KM26}\paperreference{KM26}
\bibitem[LCKM18]{LCKM18}\paperreference{LCKM18}
\bibitem[Tutte47]{Tutte47}\paperreference{Tutte47}
\bibitem[KCLL20]{KCLL20}\paperreference{KCLL20}
\bibitem[Edmonds65]{Edmonds65}\paperreference{Edmonds65}
\bibitem[CL00]{CL00}\paperreference{CL00}
\end{thebibliography}

\end{document}
